\documentclass[10pt]{article}
\usepackage[margin=0.95in]{geometry}
\usepackage{graphicx,booktabs,amsmath,amssymb,url,parskip,float}
\usepackage[hidelinks]{hyperref}

% PNAS-class macros that plain article does not define. Declared here so the
% preview compiles against exactly the same body text as the real manuscript.
\newcommand{\numberthis}{\refstepcounter{equation}\tag{\theequation}}
\newcommand{\addtabletext}[1]{\par\vspace{4pt}\footnotesize #1\par}
\newcommand{\dropcap}[1]{#1}

\title{\vspace{-2.2em}Where the neurologists are not: measuring the gap between neurological disease burden and the capacity to treat it\\[0.4em]
\large PREVIEW BUILD --- paste into the PNAS Overleaf template for the real layout}
\author{Alanna Polyak \\ \small Department of Quantitative Social Science, Dartmouth College}
\date{\today}
\begin{document}\maketitle

\begin{quote}\small\textbf{Significance statement.} Specialists per country is the statistic normally used to describe where medical capacity sits, and it answers the wrong question. Combining three World Health Organization sources across 52 countries, this study shows that the countries holding the largest share of the world's stroke burden hold the smallest share of its neurologists, by a factor of roughly thirty-four. It then shows that national counts overstate what most residents can reach, because the poorer a country is, the more completely its specialists concentrate in the capital. In low-income countries the specialist density facing a rural resident is not merely low but exactly zero, which makes the usual language of ratios inapplicable.\end{quote}

\dropcap) contains a list environment (quote, quotation, theorem, definition, enumerate, itemize...), the line after the list may have some extra indentation. If this is the case, add \parshape=0 to the end of the list environment.
\dropcap{N}eurological disorders are among the largest contributors to global
disability \cite{feigin2020}, and the workforce able to diagnose and manage them
is among the most unevenly distributed of any medical specialty. Demand is also
rising fastest where capacity is thinnest: dementia prevalence is forecast to
almost triple by 2050, with the steepest growth in low- and middle-income
regions \cite{gbd2022}. This paper asks a distributional question: are the
countries carrying the heaviest neurological disease burden the ones with the
capacity to treat it?

% Note: please start your introduction without including the word ``Introduction'' as a section heading (except for math articles in the Physical Sciences section); this heading is implied in the first paragraphs.

\subsection*{Related work}

The question is not new, and this is not the first analysis of neurological
workforce scarcity. The WHO and World Federation of Neurology Atlas has
documented that scarcity across two editions and established that resources are
severely limited in low-income countries \cite{atlas2017}. Feigin and colleagues
set out the case for translating burden evidence into workforce policy
\cite{feigin2020}, and comparative health-system statistics show specialist
supply concentrated in high-income settings, with rural residents facing
additional access barriers even inside wealthy countries \cite{oecd2023}.

What the Atlas does not do, and what this analysis adds, is link capacity to
\emph{need}. The Atlas reports medians within income groups and WHO regions,
presents no country-level workforce estimates, and includes no burden measure at
all. It can therefore say that low-income countries have fewer neurologists, but
not whether the countries with the least capacity are the ones carrying the
heaviest burden, which is the question that matters for priority-setting.
Grouped medians also suppress within-group variation entirely, so a low-income
country with functioning services and one with none are indistinguishable in the
published tables. This paper extends that literature in three ways: it joins
capacity to a burden measure at the country level, it reports a distributional
statistic rather than a central tendency, and it separates how much capacity
exists from where inside a country that capacity sits.

%The  command is used to end the formatting of the text column size that was set by the \Firstpage command. This will restore the default column size for the subsequent text.
%
%The  command should be used if the page ends in the middle of a paragraph, allowing the column formatting to continue seamlessly on the next page.

\section*{Data}

\paragraph{Sources.} Three WHO products and one survey inventory are used.
(i) The \emph{Global Dementia Observatory} \cite{gdo2017} indicator
\texttt{GDO\_q6x1\_2} gives neurologists per 100{,}000 population, self-reported
by ministries of health.
(ii) GDO indicator \texttt{GDO\_q8x3\_1} gives the furthest geographic level at
which dementia diagnostic services are available, on an ordered three-level
scale. (iii) WHO \emph{Global Health Estimates} series \texttt{SA\_0000001689}
\cite{ghe2004} gives age-standardized stroke DALYs per 100{,}000.
(iv) The WHO/WFN \emph{Neurology Atlas}, second edition \cite{atlas2017},
supplies the country to
WHO-region to World-Bank-income-group crosswalk from Annex 1 and the aggregate
workforce gradients from Tables 2--3 and Figures 11, 12 and 16. A DHS survey
inventory is used only to describe coverage for a possible subnational
extension.

\paragraph{Time window.} Capacity and service reach are 2017, the fielding year
of the second Atlas and the GDO round used. Burden is \textbf{2004}, which is
the only round of the WHO stroke DALY series exposed by the GHO API at the time
of writing. The thirteen-year offset is the single largest weakness in the
analysis and is the reason RQ1 is framed as describing a gap rather than
estimating one.

\paragraph{Unit of analysis.} In the original sources the unit differs by
product: GDO records are one row per country per indicator per round; GHE records
are one row per country per cause per age-standardization per year; the Atlas
Annex is one row per responding Member State. The unit of analysis
\emph{used here} is the \textbf{country}. Every source is collapsed to one row
per ISO3 code before merging, and no country appears twice.

\paragraph{Missingness in key fields.} The panel has 52 countries.
Neurologist density is missing for \textbf{33 of 52} (63\%), service reach for
5 (10\%), income group for 13 (25\%), and stroke burden for none. WHO encodes
absence as the strings ``Not available'', ``Not applicable'', ``No data'' and
``Not reported'' rather than as empty cells; these are converted to missing in
\texttt{code/03\_clean\_merge.py} (lines 19--20) before any numeric coercion, so
that a placeholder string cannot silently become a zero. Missingness is not
random: countries that report a workforce count plausibly have stronger health
information systems, which biases the sample toward better-resourced countries
and therefore \emph{understates} the true gap.

\paragraph{Limitations of the sources.} All WHO instruments here are ministry
questionnaires, self-reported and unaudited, with no full-time-equivalent
adjustment, so a neurologist working one clinic day a week counts the same as
one working five. Stroke stands in for neurological burden generally. Measured
burden itself depends on diagnostic capacity, which is the exposure of interest,
so the two are entangled by construction; this is addressed by bounding rather
than by adjustment.

\section*{Methods}

\paragraph{Cleaning.} \texttt{code/03\_clean\_merge.py} performs every cleaning
step. WHO placeholder strings are mapped to missing (line 20); the ordered
service-reach question is mapped from text to an integer scale via
\texttt{ACC\_SCALE} in \texttt{utils.py} (line 36); and two binary indicators are
derived from that scale, \texttt{beyond\_capital} for levels $\ge 2$ and
\texttt{reaches\_rural} for level $3$ (lines 42--43).

\paragraph{Merging.} Capacity and service reach are combined with an
\textbf{outer join on ISO3}, deliberately, so that a country reporting one but
not the other is retained rather than dropped; this yields 52 rows. Burden and
the Atlas crosswalk are then attached with \textbf{left joins on ISO3},
preserving those 52 rows. All matching is \textbf{exact on ISO3 country code},
not fuzzy on country name, which avoids the standard failure mode where
``Republic of Korea'' and ``Korea, Rep.'' fail to match. The script prints row
counts before and after every join, and the printed counts appear in
\texttt{output/log\_03\_clean\_merge.txt}.

\paragraph{Variable definitions.} The \textbf{burden-to-capacity share ratio}
for country $i$ is
%
\begin{align*}
R_i = \frac{b_i / \sum_j b_j}{c_i / \sum_j c_j} \numberthis \label{eqn:ratio}
\end{align*}
%
where $b_i$ is the age-standardized stroke DALY rate and $c_i$ is neurologist
density. Both terms are shares, so $R_i$ is invariant to the units of either
input. $R_i = 1$ means country $i$ holds exactly as large a share of the
sample's burden as of its workforce; $R_i = 13.1$ means the same staff are
stretched over thirteen times the need. It is undefined at $c_i = 0$, and those
countries are reported separately rather than dropped.

The \textbf{Urban Concentration Index} and \textbf{effective rural density} are
%
\begin{align*}
\mathrm{UCI}_g &= \frac{p^{\text{cap}}_g - p^{\text{rur}}_g}{100} \numberthis \label{eqn:uci} \\
D^{\text{eff}}_g &= m_g \cdot \frac{p^{\text{rur}}_g}{100} \numberthis \label{eqn:erd}
\end{align*}
%
where $p^{\text{cap}}_g$ and $p^{\text{rur}}_g$ are the percentages of countries
in income group $g$ reporting neurologists in the capital and in rural areas,
and $m_g$ is the group's median density. UCI is bounded $[0,1]$; a value of 1
means every country has capital coverage and none has rural coverage.
$D^{\text{eff}}_g$ rescales national density by the share of countries where a
rural resident could reach anyone at all.

Inequality is measured by the \textbf{Gini coefficient} of neurologist density,
computed by the sorted-index formula in \texttt{utils.py} (line 214) and
independently by trapezoidal integration under the Lorenz curve (line 233). The
two agree to $1.1\times10^{-16}$, which is a check on the implementation rather
than a result.

\paragraph{Inference, and why no country-level regression is reported.} This is
a \textbf{descriptive} analysis. There is no source of causal identification: no
randomization, no instrument, no policy discontinuity, and national income drives
both capacity and measured burden. No causal effect is estimated or implied.

With 19 to 35 observations, a fitted country-level slope would be badly
underpowered, so no slope is reported as a headline result. Countries are binned
by World Bank income group and the groups compared directly, with a Welch
$t$-test, a Mann-Whitney $U$, a 20{,}000-draw permutation test that assumes
nothing about distributional shape, a one-way ANOVA and Kruskal-Wallis across
three bins, a chi-square on the reach-by-income contingency table, and a
logistic regression of rural reach on high income. A regression \emph{is} fitted
in \texttt{code/07\_regression\_and\_interactions.py}, but only to demonstrate
two solvers agreeing (normal equations against gradient descent, to
$4.4\times10^{-16}$) and to test an interaction the sample cannot resolve.

\section*{Results}

\paragraph{The analysis sample.} Table~\ref{tab:sample} characterizes it. Of 52
countries in the merged panel, 19 support the RQ1 comparison and 35 the RQ3
comparison; the difference is entirely driven by which WHO indicators a country
chose to answer.

\begin{table}[t!]
\centering
\caption{The analysis sample. Row counts are after the merges described in
Methods; every source is 2017 except stroke burden, which is 2004. The two
analysis samples differ because reporting a workforce count and reporting a
service-reach level are separate questionnaire items, answered by overlapping
but non-identical sets of countries.}
\label{tab:sample}
\begin{tabular}{lr}
\toprule
Measure & $n$ \\
\midrule
Countries in merged panel & 52 \\
\quad with neurologist density (WHO GDO, 2017) & 19 \\
\quad with service-reach level (WHO GDO, 2017) & 47 \\
\quad with stroke DALY rate (WHO GHE, 2004) & 52 \\
\quad with World Bank income group (Atlas, 2017) & 39 \\
\midrule
RQ1 analysis sample: capacity and burden & 19 \\
RQ3 analysis sample: reach and income group & 35 \\
\bottomrule
\end{tabular}

\addtabletext{Countries reporting zero neurologists are retained in the panel but
cannot be placed on the ratio scale of Eq.~\ref{eqn:ratio}.}
\end{table}

\subsection*{Capacity is not aligned with burden}

On the share ratio, the median high-income country sits at \textbf{0.38} and the
median low- or middle-income country at \textbf{13.12}, a factor of
\textbf{34.3}. Myanmar reaches 69.7, Bangladesh 49.1, Togo 43.9. At the other
end Switzerland sits at 0.07. Eswatini and Fiji report \emph{zero} neurologists
and cannot be placed on the scale at all, while carrying stroke DALY rates of
1{,}536 and 994 per 100{,}000. Figure~\ref{fig:ratio} shows the full ordering.

\begin{figure}%[tbhp]
\centering
\includegraphics[width=0.74\linewidth]{figures/f01_alignment_ratio.png}
\caption{Burden-to-capacity share ratio by country, logarithmic axis. A ratio
above 1.0 means the country holds a larger share of the sample's stroke burden
than of its neurologists; below 1.0 the reverse. The axis is logarithmic because
the ratio spans three orders of magnitude, which also places parity at the
center. Capacity is WHO Global Dementia Observatory neurologist density per
100{,}000 (2017); burden is WHO Global Health Estimates age-standardized stroke
DALYs per 100{,}000 (2004). $n = 17$ countries placeable on the scale; Eswatini
and Fiji report zero neurologists and are excluded from the axis but discussed in
the text. Median high-income 0.38, median low/middle-income 13.12, a factor of
34.3.}
\label{fig:ratio}
\end{figure}

\subsection*{Capacity is distributed very unequally}

The Gini coefficient of neurologist density across the 19 reporting countries is
\textbf{0.61}, with a 95\% bootstrap interval of \textbf{[0.43, 0.73]}: high
inequality, imprecisely estimated. The poorer half of countries holds about
\textbf{6\%} of the sampled workforce. Inequality is greater between income
groups than within them, with a Gini of 0.27 within high-income countries and
0.49 within low- and middle-income countries. The focal contrast is Switzerland
at 13.10 per 100{,}000 against Myanmar at 0.07, a \textbf{187-fold} difference.
Dropping any single country moves the Gini by at most 0.06, so no one reporter
carries the result.

\begin{figure}%[tbhp]
\centering
\includegraphics[width=0.74\linewidth]{figures/f02_lorenz_gini.png}
\caption{Lorenz curve of neurologist density across the 19 countries reporting a
workforce count. The horizontal axis is the cumulative share of countries,
ordered from lowest density to highest; the vertical axis is the cumulative
share of the sampled workforce they hold. Perfect equality would place the curve
on the diagonal, and the shaded area between the two is half the Gini
coefficient by construction. The poorest half of countries holds 6\% of the
workforce; the Gini is 0.61 (95\% bootstrap CI 0.43--0.73). Source: WHO/WFN
Neurology Atlas, 2nd edition (2017).}
\label{fig:lorenz}
\end{figure}

\subsection*{The gap differs systematically between country types}

Table~\ref{tab:tests} reports the binned comparisons. Neurologist density
differs between income bins on every test applied, including the
distribution-free permutation test. The most striking descriptive result is that
dementia diagnostic services reach rural areas in 77\% of high-income countries
and in \textbf{none} of the four low or lower-middle income countries sampled.

\begin{table}[t!]
\centering
\caption{Binned hypothesis tests. Country-level regressions are not fitted as
headline results because 19 to 35 observations cannot support them, so countries
are binned by World Bank income group and the groups compared directly. Rank-based
and permutation tests are reported alongside the parametric ones because the
chi-square has six of nine cells with expected counts below five.}
\label{tab:tests}
\begin{tabular}{llr}
\toprule
Comparison & Test & Result \\
\midrule
Density, high vs low/middle & Welch $t$ & $t = 4.29$, $P = 0.0031$ \\
\quad same & Mann-Whitney $U$ & $P = 0.0003$ \\
\quad same & Permutation (20{,}000) & $P = 0.0003$ \\
Reach across three bins & One-way ANOVA & $F = 5.74$, $P = 0.0074$ \\
\quad same & Kruskal-Wallis & $H = 9.44$, $P = 0.0089$ \\
Reach $\times$ income bin & Chi-square & $\chi^2 = 11.81$, $P = 0.0188$ \\
Rural reach $\sim$ high income & Logistic & OR $= 7.6$, $P = 0.0098$ \\
\bottomrule
\end{tabular}

\addtabletext{All tests are two-sided. The permutation test shuffles income-group
labels 20{,}000 times and assumes nothing about distributional shape.}
\end{table}

\subsection*{The gradient steepens with rarity}

Across the 114 Atlas respondents, the high-to-low income ratio is 71-fold for
total neurological workforce, 158-fold for adult neurologists, 62-fold for
neurosurgeons and \textbf{195-fold} for child neurologists. The scarcer the
cadre, the more concentrated it is in rich countries.

\subsection*{National density overstates access}

Figure~\ref{fig:urbanrural} reports both indices. UCI falls monotonically from
0.95 in low-income countries through 0.83 and 0.78 to 0.52 in high-income ones.
Effective rural density rises across the same groups from \textbf{exactly zero}
through 0.24 and 0.59 to 3.19 per 100{,}000. The zero is not
a rounding artifact: no low-income country in the Atlas reports a single
neurologist practicing rurally. The national gap between high- and low-income
countries is 71-fold; the rural gap has no denominator at all.



\begin{figure}%[tbhp]
\centering
\includegraphics[width=0.74\linewidth]{figures/f06_urban_rural_indices.png}
\caption{National density overstates access wherever specialists stay urban.
Panel A: the Urban Concentration Index, the share of countries where
neurologists reach the capital minus the share where they reach rural areas.
Panel B: effective rural density, the national median rescaled by the share of
countries with any rural practice, which is the quantity a rural resident
actually faces. The low-income bar in Panel B is exactly zero, not rounded to
zero. Values by income group, low to high: UCI 0.95, 0.83, 0.78, 0.52;
effective rural density 0.00, 0.24, 0.59, 3.19 per 100{,}000. WHO/WFN Neurology
Atlas, 2nd edition (2017), 114 responding countries.}
\label{fig:urbanrural}
\end{figure}

\subsection*{An interaction the sample cannot resolve}

If capacity buys less alignment where services remain urban, capacity and rural
reach should interact. Only \textbf{13 countries} carry both measures. The
interaction coefficient is $-0.004$ with $P = 0.99$. That is a statement about
statistical power, not evidence that no interaction exists, and the figure title
in the repository says so explicitly rather than implying a null finding.

\subsection*{Robustness}

Three checks. A 5{,}000-draw bootstrap gives the Gini interval quoted above.
Leave-one-out moves the Gini by at most 0.06 for any single country. A bounding
exercise asks how much measurement error would be needed to overturn the central
result: burden in low- and middle-income countries would have to be
\textbf{overstated by a factor of 34}, or their capacity understated by the same
factor, for the two median ratios to level. Measurement error of that magnitude
is not plausible, and under-diagnosis in low-capacity settings pushes
\emph{measured} burden down rather than up, which widens the true gap rather
than closing it.

\section*{Discussion}

Across three metrics the finding is consistent: neurological capacity is
concentrated where measured burden is lightest, and national counts overstate
what most residents can reach. The policy implication is that specialists per
country is the wrong statistic for priority-setting. It answers a question about
supply when the operative question is about reach, and it is silent about
whether a rural resident can access anyone at all.

Six limitations bound the conclusion. The thirteen-year offset between capacity
and burden is the weakest joint. The samples are small and self-selected toward
countries with functioning health information systems, which understates the
gap. Stroke stands in for neurological burden generally. Six of nine cells in
the reach-by-income table have expected counts below five. Both WHO instruments
are self-reported and unaudited with no FTE adjustment. And the interaction test
is underpowered at 13 countries. Nothing here is a causal estimate.

Two extensions follow directly. The DHS inventory described in the repository
covers 106 countries and 443 surveys, 62 with GPS coordinates, which would
permit the same alignment question to be asked \emph{within} countries rather
than across them. And a burden series contemporaneous with the 2017 capacity
measures would remove the study's largest weakness.

\section*{Appendix A. The remaining figures}

Three figures and three tables appear in the body. The pipeline produces
fourteen figures and ten tables. The eleven held back are not weaker versions of
those shown; they answer adjacent questions, and several exist to check the
analysis rather than to say anything about neurologists. All are reproduced here
with the same self-contained captions they carry in the repository. This
appendix and the Agentic Analysis section that follows are excluded from the
page limit.

\subsection*{A.1 Extending the descriptive picture}

\begin{figure}[H]
\centering
\includegraphics[width=0.74\linewidth]{figures/f03_service_reach_by_income.png}
\caption{Dementia diagnostic service reach by World Bank income group. The
ordered scale runs capital only, capital plus main cities, reaches rural areas.
Rural reach occurs in 77\% of high-income countries and in none of the four low
or lower-middle income countries sampled. WHO Global Dementia Observatory
(2017), $n = 35$ countries with both a reach level and an income group.}
\label{fig:a-reach}
\end{figure}

\begin{figure}[H]
\centering
\includegraphics[width=0.74\linewidth]{figures/f04_atlas_workforce_by_income.png}
\caption{Median neurological workforce per 100{,}000 by income group across all
114 Atlas respondents, rather than the 19 with a usable country-level count.
The high-to-low ratio is 71-fold. Source: WHO/WFN Neurology Atlas, 2nd edition
(2017), Figure 12.}
\label{fig:a-atlasinc}
\end{figure}

\begin{figure}[H]
\centering
\includegraphics[width=0.74\linewidth]{figures/f05_atlas_where_neurologists_practice.png}
\caption{Where neurologists practice, by income group. Capital-city coverage is
near universal at every income level; rural coverage falls from 45\% to 0\%.
This is the aggregate counterpart to Fig.~\ref{fig:a-reach}, measured by a
different instrument on a larger country set, and it produces the same ordering.
WHO/WFN Neurology Atlas (2017), Figure 16, 114 responding countries.}
\label{fig:a-practice}
\end{figure}

\begin{figure}[H]
\centering
\includegraphics[width=0.74\linewidth]{figures/f10_cadre_gradient.png}
\caption{The income gradient by clinical cadre, logarithmic axis. The scarcer
the cadre, the more concentrated it is in high-income countries: 71-fold for
total workforce, 158-fold for adult neurologists, 62-fold for neurosurgeons and
195-fold for child neurologists. WHO/WFN Neurology Atlas (2017), Tables 2--3.}
\label{fig:a-cadre}
\end{figure}

\subsection*{A.2 Geography rather than economics}

\begin{figure}[H]
\centering
\includegraphics[width=0.74\linewidth]{figures/f09_regional_gradient.png}
\caption{Median workforce density by WHO region with individual reporting
countries overlaid, so within-region spread remains visible. Income group and
region are correlated but not identical, and the overlay is what shows the
difference. WHO/WFN Neurology Atlas (2017) and WHO Global Dementia Observatory
(2017).}
\label{fig:a-region}
\end{figure}

\subsection*{A.3 Paired instrument checks}

Each panel places an Atlas aggregate beside its country-level counterpart from
the Global Dementia Observatory. Two instruments collected independently, by
different teams, with different questions, ordering countries the same way is
weak evidence that neither is simply wrong. It is not independent confirmation,
because both are ministry self-reports.

\begin{figure}[H]
\centering
\includegraphics[width=0.74\linewidth]{figures/f11_paired_workforce.png}
\caption{The same workforce gradient measured two ways. Panel A reports the
median within each income group and cannot identify individual countries; Panel
B reports each country separately but covers far fewer of them. They disagree on
level, because Panel A counts three cadres and Panel B counts adult
neurologists only, and agree on direction and steepness.}
\label{fig:a-pairwf}
\end{figure}

\begin{figure}[H]
\centering
\includegraphics[width=0.74\linewidth]{figures/f12_paired_rural_access.png}
\caption{The same rural-access gradient measured two ways, with the Atlas
aggregate beside the country-level Global Dementia Observatory extract.}
\label{fig:a-pairrural}
\end{figure}

\subsection*{A.4 Groundwork and verification}

\begin{figure}[H]
\centering
\includegraphics[width=0.74\linewidth]{figures/f13_dhs_coverage.png}
\caption{Demographic and Health Survey coverage across the sampled countries:
3{,}018 datasets, 106 countries, 443 surveys, 62 with GPS coordinates. This is
groundwork for asking the alignment question \emph{within} countries rather than
across them. It is not a result.}
\label{fig:a-dhs}
\end{figure}

\begin{figure}[H]
\centering
\includegraphics[width=0.74\linewidth]{figures/f07_regression_solvers.png}
\caption{The two estimators of the same regression, closed-form normal equations
against gradient descent, agreeing to $4.4 \times 10^{-16}$. This verifies the
implementation and says nothing about neurologists. It is included because a
hand-written estimator that agrees with \texttt{statsmodels} to machine
precision is the cheapest available evidence that the algebra is right.}
\label{fig:a-solvers}
\end{figure}

\begin{figure}[H]
\centering
\includegraphics[width=0.74\linewidth]{figures/f08_interactions.png}
\caption{The interaction the 13-country sample cannot resolve, beside the
comparison it can. Panel A: two nearly parallel fits, interaction coefficient
$-0.004$ at $P = 0.99$. Panel B: rural reach is far more common in high-income
countries. The title states what the figure shows, including that the sample
cannot resolve the first question.}
\label{fig:a-inter}
\end{figure}

\begin{figure}[H]
\centering
\includegraphics[width=0.74\linewidth]{figures/f14_robustness.png}
\caption{Robustness. Bootstrap distribution of the Gini, leave-one-out
sensitivity, and the bounding exercise. Dropping any single country moves the
Gini by at most 0.06, and burden would have to be overstated by a factor of 34
in low- and middle-income countries for the central result to reverse.}
\label{fig:a-robust}
\end{figure}

\subsection*{A.5 Omitted tables}

Seven tables hold the numbers behind results stated in prose above and are not
reproduced here: the Atlas gradients (\texttt{t03}), the regression output
(\texttt{t05}), the interaction tests (\texttt{t06}), the regional comparison
(\texttt{t07}), the paired-instrument correlations (\texttt{t08}), DHS coverage
counts (\texttt{t09}), and the leave-one-out and bounding results
(\texttt{t10}). All are committed as CSV in \texttt{output/tables/}.

\section*{Agentic Analysis}

An AI assistant (Claude, Anthropic) was used throughout. Transcripts are in
\texttt{docs/agentic/}. This section is excluded from the page limit.

My own role is quickly stated. The research questions, the choice of a share
ratio over a difference, the decision to bin rather than regress, and the
bounding exercise are mine; I executed and read every line of code, and traced
every reported number back to a source table. The space below goes to the four
episodes where the assistant was wrong, because the diagnoses transfer and the
successes do not.

\paragraph{Rejected: a regression presented as a headline.} The assistant
proposed fitting the alignment ratio on capacity and income and reporting the
slope as the paper's main quantitative result, with a standard error and a
$p$-value.

\emph{How it was wrong:} the estimate is not identified in any useful sense on
19 observations, but the output does not look that way. A coefficient with a
standard error carries an implicit claim to precision that the sample cannot
support, and nothing in the fitted object flags this. I replaced it with binned
comparisons and a 20{,}000-draw permutation test, which make a weaker claim
honestly. The regression survives only in Fig.~\ref{fig:a-solvers}, as a check
that two solvers agree.

\paragraph{Rejected: dropping the countries that break the plot.} Eswatini and
Fiji report zero neurologists, so the share ratio is undefined for them and they
cannot be drawn on a logarithmic axis. The assistant proposed excluding them.

\emph{How it was wrong:} the suggestion optimised for a clean figure over a
correct sample. Those two countries are the most extreme instances of the
phenomenon the paper is about, and their stroke DALY rates, 1{,}536 and 994 per
100{,}000, are among the highest in the sample. Dropping them would have removed
the strongest evidence for the claim while improving the axis. They are reported
separately in the text and named in the caption instead.

\paragraph{Failure: a figure title asserting a finding the data did not support.}
An earlier revision produced a panel titled ``adding capacity does not help
equally in every country,'' describing an interaction whose coefficient was
$-0.004$ at $P = 0.99$.

\emph{How it was wrong:} the code ran, the figure rendered, and the title stated
a substantive conclusion drawn from a null result on 13 observations. Nothing in
the pipeline connects a title string to the significance of what it describes,
so this class of error is invisible to testing. The panel now reads ``the
interaction the 13-country sample cannot resolve, beside the comparison it
can,'' which states what the figure shows including the part it cannot answer.

\paragraph{Failure: silent truncation in the PDF parser, and an over-correction.}
The Neurology Atlas wraps two long country names across lines. The first parser
took only the first line, yielding ``The former Yugoslav Republic of'' as a
country name.

\emph{How it was wrong, first attempt:} the parse succeeded, produced 133 rows
as expected, and the corrupted values sat in a column nobody reads closely. Row
counts are the usual check and row counts were correct.

\emph{How it was wrong, second attempt:} the assistant's fix appended any short
following line to the country name. This over-corrected, gluing contributor
names onto countries and producing entries such as ``Austria Grisold, Regina
Katzenschlager.'' The rule that held keys on the observation that a wrapped
country name is cut mid-phrase and ends in a function word, which no contributor
name does. Eight assertions against values read off the page now guard the
parse.

\paragraph{Failure: a data exposure I was not warned about.} The DHS inventory
the assistant built retained the authenticated download URLs in a committed
column. Those URLs are tied to my approved DHS account.

\emph{How it was wrong:} nothing downstream used the column, so its presence
changed no result and produced no error. It would have been published to a
public repository. \texttt{02\_build\_dhs\_inventory.py} now discards the URLs
at parse time, and the exclusion is stated in the script docstring so it cannot
be reintroduced by someone rebuilding the file.

\paragraph{What this implies.} All four share a shape: output that is internally
consistent and externally wrong, produced with the same confidence whether right
or not. None would have been caught by a test suite, because in each case the
code did what it said. Each surfaced only when I asked what a quantity
\emph{meant} rather than whether the script ran. AI assistance moved the
bottleneck from writing code to verifying that the code answered the intended
question, and the second is the harder job.

\section*{Materials and Methods}

All code, committed data and figures are at
\url{https://github.com/alannap27/qss20-neuro-access}, organized as numbered
scripts in \texttt{code/}, committed inputs in \texttt{data/}, and figures and
tables in \texttt{output/}. Shared functions are in \texttt{code/utils.py} and
figure layout rules in \texttt{code/figstyle.py}. No script contains a hardcoded
absolute path. The WHO inputs are public and committed; the Neurology Atlas PDF
and the DHS manifest are excluded, the first because it is a copyrighted
publication and the second because its URLs are tied to an approved DHS account.
Analyses used Python 3.10 with pandas, NumPy, SciPy, statsmodels and Matplotlib;
versions are pinned in \texttt{requirements.txt}. Random seeds are fixed
throughout.

\bibliographystyle{plain}
\bibliography{qss20}
\end{document}
